\section{Cards M: Memory Hierarchy}\label{grp:mem}
$\AMAT=t_{hit}+m\times P$;\quad $\CPI=\CPI_{base}+\text{refs/instr}\times m\times P$ (slide: $2+1.2\times0.03\times64=4.3$). Levers: cut \textbf{penalty} (M1, M8, M9, M5), \textbf{miss rate} (M2--M7, M13, M14), \textbf{both via parallelism} (M10--M12), \textbf{hit time} (M15, M16). Hit time matters most: it sets the clock for \emph{every} instruction. \textbf{First diagnose the miss type:}

\small
\begin{tabularx}{\textwidth}{L{2.4cm}L{5.4cm}L{4.4cm}X}
\toprule
\tblhead{\textbf{Miss type} & \textbf{Meaning / test} & \textbf{Fixed by} & \textbf{NOT fixed by}}\\
\midrule
Compulsory & First touch; happens even in an infinite cache & Larger blocks, prefetch & Bigger cache, associativity \\
Capacity & Working set $>$ cache; happens even fully associative & Bigger cache, blocking, fusion & Associativity \\
Conflict & Too many blocks map to one set; vanishes with full associativity & Associativity, victim cache, padding, pseudo-assoc. & Compulsory/capacity fixes \\
Coherence & Line invalidated by another PE's write (true or false sharing) & Privatise, pad/align, smaller blocks & Bigger caches \\
\bottomrule
\end{tabularx}
\normalsize
\begin{multicols}{2}
\card{M1}{Multilevel caches (L2/L3)}{Static}
\fx{Large miss penalty to DRAM (``best penalty reduction'').}
\fh{L2 catches L1 capacity/conflict misses. $\AMAT=t_{L1}+m_{L1}(t_{L2}+m^{loc}_{L2}P)$. Judge L2 by \emph{global} miss rate $m_{L1}\times m^{loc}_{L2}$; L2 size matters most.}
\fu{Fast core, far memory (desktop/server: L1+L2+L3; PMD: L1 64\,KB + L2 256\,KB, no L3).}
\fb{Bigger L1: keeps L1 fast (clock) while catching misses.}
\fq{1000 refs, 40 L1 and 20 L2 misses: $m_{L1}=4\%$, local L2 50\%, global 2\%; $\AMAT=1+0.04(10+0.5\times100)=3.4$; stalls/instr $=(3.4-1)\times1.5=3.6$.}
\fj{L1 misses cost [100] cycles; an L2 of [size] with [x\%] local miss rate cuts AMAT from [a] to [b] without slowing the L1 hit.}

\card{M2}{Larger block size}{Static}
\fx{Compulsory misses (spatial locality).}
\fu{Sequential/streaming data, large L2/L3 (64--128\,B).}
\fa{Small L1 with 128--256\,B blocks; random access; heavy sharing.}
\fw{Raises miss penalty; in small caches raises conflict/capacity misses; more false sharing.}
\fq{AMAT-optimal block: 4\,KB cache 32\,B; 16--256\,KB caches 64\,B (model: 80\,cc + 2\,cc per 16\,B). 256\,B only wins at 256\,KB.}
\fj{The access is [sequential], so moving from [32] to 64\,B blocks halves compulsory misses; going to 256\,B would raise the penalty and conflicts in a [16\,KB] L1 (AMAT 3.32 at 64\,B vs.\ 4.69 at 256\,B).}

\card{M3}{Larger cache}{Static}
\fx{Capacity misses.}
\fu{Working set slightly above current size, mainly at L2/L3.}
\fb{Associativity for capacity misses (associativity does not fix capacity).}
\fa{Working set far larger (use blocking \cp{M14}) or streaming.}
\fw{Longer hit time, area, power; useless for streaming (no reuse).}
\fj{The working set [x\,MB] just exceeds the [y\,MB] L2, so these are capacity misses; enlarging the off-critical-path L2 removes them without slowing the L1.}

\card{M4}{Higher associativity}{Static}
\fx{Conflict misses.}
\fh{$n$ ways per set. 8-way $\approx$ fully associative. 2:1 rule: 2-way of size $N/2\approx$ DM of size $N$ ($<$128\,KB).}
\fu{\textbf{At L2/L3} (extra hit time $\ll$ penalty); tiny L1 (4--8\,KB); many threads sharing a cache.}
\fa{Clock-critical L1 $\ge$16\,KB: DM has lowest AMAT once the clock stretch is charged.}
\fw{Hit time $\Rightarrow$ clock $\times1.10$ (2-way), $\times1.12$ (4), $\times1.14$ (8).}
\fq{L2: DM $10+0.25\times100=35$ vs.\ 2-way $10+0.20\times100=30$. Lab 3: L1 associativity fixed a conflict storm but L2/L3 miss \emph{rates} rose while AMAT improved: judge by AMAT.}
\fj{Addresses [a] and [b] map to the same set (power-of-2 stride), so these are conflict misses; making the L2 8-way removes them at a negligible hit-time cost, whereas a more associative L1 would stretch the clock by 10--14\%.}

\card{M5}{Victim cache}{Dynamic}
\fx{Conflict misses of a fast direct-mapped L1.}
\fh{Small fully-associative buffer of just-evicted blocks, probed with the miss; swap on hit.}
\fb{Raising L1 associativity: keeps the DM hit time (clock) and recovers much of 2-way's miss rate.}
\fa{Capacity misses or streaming (victims are never reused).}
\fq{Removes 20--95\% of misses of a 4\,KB DM cache with 1--5 entries; Athlon 8 entries. Lab 4: 61.6\% fewer cycles on a conflict pattern, nothing on the full pipeline; cutoff at ($N-$ways) entries.}
\fj{The DM L1 ping-pongs between [arrays a and b]; a 4-entry victim cache catches the just-evicted block, keeping the 1-cycle DM hit, which 2-way associativity (clock $\times1.10$) would not.}

\card{M6}{Way prediction}{Dynamic}
\fx{Hit time and energy of set-associative caches.}
\fh{Predict the way; set the mux early; compare one tag first; mispredict $\Rightarrow$ extra cycle.}
\fu{PMDs, I-caches (sequential, $>$85\% accurate on Alpha 21264 2-way).}
\fb{Full parallel compare: one way active $\Rightarrow$ less energy, near-DM hit time.}
\fw{Predictor bits; an extra cycle on a wrong way.}
\fj{In a battery-limited core, way prediction compares one tag instead of [n], keeping 2-way miss rate at near-DM hit time and energy.}

\card{M7}{Pseudo-associative (column-associative)}{Dynamic}
\fx{Conflict misses at DM speed.}
\fh{DM probe; on miss probe a second slot (flip top index bit) sequentially; swap frequent block to fast slot.}
\fb{2-way at L2: DM-speed primary hit with near 2-way miss rate.}
\fw{Variable 1-or-2-cycle hit: bad for L1 pipeline $\Rightarrow$ use at \textbf{L2} (MIPS R10000).}
\fq{AMAT 2\,KB: DM 5.90, 2-way 4.90, pseudo 4.95; 128\,KB: 1.50 / 1.45 / 1.371.}

\card{M8}{Critical word first / early restart}{Dynamic}
\fx{Waiting for a whole large block.}
\fh{Send the requested word first (or restart as soon as it arrives).}
\fu{Large blocks; weak if the next access is the next word anyway.}
\fw{Little gain with small blocks or when the next access is the next word.}

\card{M9}{Read priority over writes; merging write buffer}{Dynamic}
\fx{Read misses waiting behind buffered writes; buffer-full stalls.}
\fh{On a read miss check the write buffer for the address and go ahead if no match; write-back: buffer the dirty victim, read first. Merge writes to the same block into one entry.}
\fu{Write-through or write-heavy code.}
\fb{Waiting for the write buffer to drain before every read miss.}
\fq{Waiting for a 4-word buffer to drain: $\approx$50\% worse read-miss penalty (MIPS 1000).}
\fj{Write-through traffic [given] fills the buffer; letting reads bypass non-matching writes and merging same-block writes removes the read-miss delay and buffer-full stalls.}

\card{M10}{Non-blocking cache (hit-under-miss, MSHRs)}{Dynamic}
\fx{CPU idle while one miss is served.}
\fh{Keep serving hits during a miss; with several MSHRs, overlap several misses (memory-level parallelism).}
\fu{OoO cores, independent misses (several streams, many threads).}
\fb{Associativity for cutting stall time: no hit-time increase.}
\fa{Dependent misses (pointer chase): MLP $=1$. Lab 4: no gain on a chase, 8$\times$ on 8 independent streams, plateau at MSHR $=8$.}
\fn{\textbf{OoO CPU} (an in-order blocking CPU stalls anyway) and \textbf{banked memory} \cp{M18}.}
\fq{OoO hides L1-miss/L2-hit, not most L2 misses.}
\fj{The OoO core has [k] independent loads missing together; with [8] MSHRs their misses overlap, so the effective penalty falls by up to [k]$\times$ without changing hit time.}

\card{M11}{Hardware prefetch (stream buffers, stride)}{Dynamic}
\fx{Compulsory/capacity misses of regular access.}
\fh{Detect sequential/strided pattern; fetch ahead into a buffer.}
\fu{Streaming, sequential, constant stride; spare bandwidth.}
\fb{Compiler prefetch: no instruction overhead, adapts to the run-time pattern.}
\fw{Useless/harmful when untimely, polluting, or bandwidth-starved (Lab 4: misses 57.4\%$\rightarrow$23.2\% yet slower).}
\fn{Spare memory bandwidth; non-blocking cache; throttling.}
\fq{I-stream buffer 1/4/16 blocks catch 15--25/50/72\% of I-misses; 1 data stream $-25\%$, 4 streams $-43\%$; prefetch AMAT $2+0.011\times0.25\times1+0.011\times0.75\times50=2.415$.}
\fj{The [stream] is sequential and bandwidth is [x\% used], so a stream prefetcher fetches ahead and hides the [n]-cycle latency; I would disable it if bandwidth were saturated.}

\card{M12}{Compiler-controlled prefetch}{Static}
\fx{Misses in loops whose stride/distance is known.}
\fh{Non-faulting prefetch instructions ahead of use (\texttt{prefetch(b[j+7])}).}
\fb{HW prefetch when the pattern is known at compile time or irregular-but-computable.}
\fa{Irregular patterns; bandwidth saturated.}
\fw{Instruction overhead; wrong distance $\Rightarrow$ too early (evicted) or too late.}
\fn{Non-blocking cache; spare issue slots (benefit must exceed issue cost).}

\card{M13}{Compiler layout: interchange, fusion, merging, padding, SoA}{Static}
\fx{Poor locality and conflicts, at zero hardware cost.}
\fh{\emph{Interchange}: inner loop walks memory sequentially. \emph{Fusion}: reuse data while cached. \emph{Merging arrays} (AoS) for fields used together; \emph{struct-of-arrays} for SIMD/GPU. \emph{Padding}: break power-of-2 aliasing and false sharing. Procedure reordering for I-cache.}
\fu{Array/loop code (scientific, image), embedded designs where no hardware can be added.}
\fb{Any hardware fix: zero area and power.}
\fa{Pointer structures; transformations illegal due to dependences.}
\fq{Interchange $3{,}145{,}728\rightarrow196{,}608$ misses ($-93.75\%$); fusion $262{,}144\rightarrow196{,}608$ ($-25\%$); McFarling code reordering $-50$/$-75\%$ (2/8\,KB I-cache).}
\fj{The inner loop walks a column of a row-major array [given], one useful word per 64\,B line; interchanging the loops makes it sequential (slide: $-93.75\%$ misses) at no hardware cost.}

\card{M14}{Blocking (tiling)}{Static}
\fx{Capacity misses when rows \emph{and} columns are reused (matrix multiply).}
\fh{Work on $B\times B$ tiles; three tiles must fit in cache.}
\fu{Matrix multiply, stencils, any reuse along two dimensions.}
\fb{Interchange (cannot help both access directions); bigger cache (free).}
\fw{Tile-size tuning, more complex code.}
\fq{Misses $2N^3+N^2\rightarrow2N^3/B+N^2$; $N=1024$, 32\,KB, $B=32$: $\approx1.07$\,billion $\rightarrow\approx2.1$\,million (simplified model).}
\fj{The $[N]\times[N]$ matrices exceed the cache, so every column pass misses; tiles of $B=[32]$ (three tiles $=[12]$\,KB fit the [32]\,KB L1) cut misses from $2N^3$ to $2N^3/B$.}

\card{M15}{Small, simple L1; pipelined access; virtual indexing}{Static}
\fx{Hit time limiting the clock.}
\fh{Small DM/low-way L1; pipeline write tag check with previous data write; index the cache in parallel with TLB translation (index within page offset; or PIDs + anti-aliasing/page colouring for virtual caches).}
\fu{Every high-clock design.}
\fb{Bigger associative L1: from 16\,KB up, DM has the lowest AMAT once the clock is charged.}
\fw{More misses (offset by L2).}
\fq{PIDs in tags save 0.6--4.3\% miss rate vs.\ flushing on a switch.}
\fj{The L1 hit sets the [x\,GHz] clock, so I keep it small and low-associativity and move capacity to L2.}

\card{M16}{Trace cache}{Dynamic}
\fx{Fetch bandwidth lost at taken branches (one every 5--10 instructions).}
\fh{Store dynamic instruction traces across branches (Pentium 4).}
\fu{Wide-issue front ends on branchy code with frequent taken branches.}
\fa{Energy-constrained designs.}
\fw{Complex mapping, duplicated instructions, area/energy.}

\card{M17}{Wider main memory / bus}{Static}
\fx{Block transfer bandwidth.}
\fh{$N$-word-wide memory and bus.}
\fu{The bus can be widened (pins, cost acceptable); small blocks.}
\fb{Interleaving at 2-word blocks (3.54 vs.\ 3.63 CPI).}
\fw{Cost; mux on critical path (place between L1 and L2); ECC needs read-modify-write; expansion in big increments.}

\card{M18}{Interleaved banks / independent banks / more channels}{Static}
\fx{Bandwidth for sequential blocks, multiple misses, many cores/lanes.}
\fh{Bank $=$ address mod banks; accesses overlap on a shared bus; independent banks/channels serve independent requests.}
\fu{Streaming, vector/SIMD/GPU, non-blocking caches, multicore.}
\fb{Wide memory when the bus cannot be widened (pins/cost); best with larger blocks.}
\fw{Many chips, harder expansion; controller complexity for independent banks.}
\fq{Need banks $\ge$ access time in cycles. Slide CPI: 1-word 4.30; 2-word simple 5.07 / interleaved 3.63 / wide 3.54; 4-word 5.69 / \textbf{3.09} / 3.84.}
\fj{[k] cores/lanes stream [x\,GB/s] but one channel gives [y]; adding channels/banks lets their misses proceed in parallel, which a larger cache cannot do for data used once.}

\card{M19}{Virtual memory and TLB}{Dynamic}
\fx{Protection, many processes, memory larger than DRAM.}
\fh{Pages (4--8\,KB, no external fragmentation) mapped by a page table; TLB caches translations (128--256 entries, fully associative on high-end).}
\fb{Segmentation: fixed pages avoid external fragmentation and need one-word addresses.}
\fw{TLB misses when the working set exceeds TLB reach (larger pages\bs).}
\fq{32-bit VA $=$ 20-bit VPN + 12-bit offset.}


\card{M20}{Write policy: write-back vs.\ write-through}{Static}
\fx{Write traffic to the next level.}
\fh{\emph{Write-through}: every write also goes to the next level (via a write buffer); cache always clean. \emph{Write-back}: write only the cache, mark the block dirty, write it out on eviction.}
\fu{Write-back: write-heavy code, multicore (less bus traffic), L2/L3. Write-through: simple L1 backed by an L2; easy recovery and coherence.}
\fb{Write-back cuts traffic when a block is written many times; write-through keeps the lower level current.}
\fw{Write-back: dirty evictions lengthen misses (buffer them, \cp{M9}). Write-through: bandwidth on every write; needs a write buffer.}
\fn{Write-through: merging write buffer + read priority \cp{M9}.}
\fj{The loop rewrites the same [block] [k] times; write-back sends it to memory once on eviction instead of [k] times, freeing bus bandwidth for the other cores.}
\end{multicols}

\small
\begin{tabularx}{\textwidth}{L{1.4cm}L{1.3cm}L{1.3cm}L{1.3cm}L{1.3cm}X|L{1.2cm}L{1.1cm}L{1.1cm}L{1.1cm}L{1.1cm}}
\toprule
\tblhead{\textbf{AMAT} & \textbf{1-way} & \textbf{2-way} & \textbf{4-way} & \textbf{8-way} & & \textbf{Block} & \textbf{4K} & \textbf{16K} & \textbf{64K} & \textbf{256K}}\\
\midrule
4\,KB & 3.44 & 3.25 & \textbf{3.22} & 3.28 & & 16\,B & 8.03 & 4.23 & 2.67 & 1.89 \\
8\,KB & 2.69 & 2.58 & \textbf{2.55} & 2.62 & & 32\,B & \textbf{7.08} & 3.41 & 2.13 & 1.59 \\
16\,KB & \textbf{2.23} & 2.40 & 2.46 & 2.53 & & 64\,B & 7.16 & \textbf{3.32} & \textbf{1.93} & \textbf{1.45} \\
64\,KB & \textbf{1.92} & 2.14 & 2.18 & 2.25 & & 128\,B & 8.47 & 3.66 & 1.98 & 1.47 \\
256\,KB & \textbf{1.32} & 1.66 & 1.74 & 1.82 & & 256\,B & 11.65 & 4.69 & 2.29 & 1.55 \\
\bottomrule
\end{tabularx}
\normalsize
Left: AMAT by size and associativity \emph{with clock stretch charged} (DM wins from 16\,KB). Right: AMAT by block size and cache size. Hierarchy by class: \textbf{PMD} L1 64\,KB, L2 256\,KB, 1--2\,GB DRAM, flash; \textbf{desktop} + L3 8--32\,MB, 4--64\,GB; \textbf{server} L3 16--64\,MB, 32--256\,GB DRAM, flash and disk.


\small
\begin{tabularx}{\textwidth}{L{4.4cm}C{1.3cm}C{1.3cm}C{1.3cm}C{1.3cm}X}
\toprule
\tblhead{\textbf{Technique (card)} & \textbf{Penalty} & \textbf{Miss rate} & \textbf{Hit time} & \textbf{HW cost} & \textbf{Used in / comment}}\\
\midrule
Multilevel caches \cp{M1} & $+$ & & & 2 & Costly, universal \\
Critical word first \cp{M8} & $+$ & & & 2 & Widely used \\
Read priority / merging buffer \cp{M9} & $+$ & & & 1 & Write-through \\
Victim cache \cp{M5} & $+$ & $+$ & & 2 & AMD Athlon (8 entries) \\
Larger block \cp{M2} & $-$ & $+$ & & 0 & Pentium 4 L2 128\,B \\
Larger cache \cp{M3} & & $+$ & $-$ & 1 & Esp.\ L2 \\
Higher associativity \cp{M4} & & $+$ & $-$ & 1 & Widely used \\
Way prediction \cp{M6} & & & $+$ & 2 & UltraSPARC III I-cache \\
Pseudo-associative \cp{M7} & & $+$ & & 2 & MIPS R10000 L2 \\
Compiler layout/blocking \cp{M13} \cp{M14} & & $+$ & & 0 & Software effort \\
Non-blocking \cp{M10} & $+$ & & & 3 & All OoO CPUs \\
HW / compiler prefetch \cp{M11} \cp{M12} & $+$ & $+$ & & 2--3 & Needs bandwidth / non-blocking \\
Small simple L1, pipelined, no translation \cp{M15} & & ($-$) & $+$ & 0--2 & Alpha 21164 \\
Trace cache \cp{M16} & & & $+$ & 3 & Pentium 4 \\
\bottomrule
\end{tabularx}
\normalsize
($+$ improves, $-$ hurts; HW 0 none, 3 highest.) Conflicts of interest: larger blocks hurt penalty; size/associativity hurt hit time; small/simple hurts miss rate.
